\section{Was aus den Quellenantworten im affinen Netz entsteht}
Unter der bestehenden unitären $E_{8,1}$-Vakuumprämisse liefert
\[
 i_n(a)=\frac1{\sqrt n}J_{-n}(a)\ket\Omega
\]
eine normtreue Einbettung. Die adjungierten Nullmoden handeln richtig, und $[L_0,J_{-n}]=nJ_{-n}$. Mit der zusätzlichen Randzeitwahl
\[
 H_{\mathrm{Rand}}=\frac{2\pi v}{L}\left(L_0-\frac c{24}\right),
 \qquad c=8,
\]
entsteht eine lineare Antwort bei $\omega=vq_n$, $q_n=2\pi n/L$. Für die quartischen Richtungen lautet das Gewicht
\[
 S_{lr}^{(n)}(d\omega)=\frac nL
 \langle\psi_r,\psi_l\rangle^3\delta_{vq_n}(d\omega).
\]
Dies ist ein präziser positiver Anschluss. Die Randlänge $L$, die Geschwindigkeit $v$ und die Identifikation mit der Rohquelle werden dadurch nicht hergeleitet. \src{20}

Die 20 quartischen Richtungen schließen mit ihren Adjungierten nur $A_8$. Ein bereits vorhandener linearer Begleiter in $\Lambda^0(\C^5)\otimes4$ vervollständigt über tatsächliche Chevalleyklammern ganz E8. Ein 120er-Frame in einem 24-dimensionalen Stromraum besitzt im gewählten kubisch-plus-linearen Ansatz das Gewichtsverhältnis $5:1$. Weder die Dimension 24 noch dieses Verhältnis werden als neue Teilchenzählung ausgegeben. \src{20}

\section{Das alte Bell10 erscheint aus echten Stromprodukten wieder}
Die tatsächlichen Grad-2-Produkte der 24 Ströme erzeugen einen 50-dimensionalen Raum. Für die 60 Produktvektoren $r_l$ gilt
\[
 \mathcal G=\sum_l\ket{r_l}\bra{r_l}
 =\frac34I_{50}+\frac94P_{10}.
\]
Der alte Bell10-Raum ist damit der eindeutige oberste \emph{Antwortbereich}: zehn Eigenwerte sind $3$, vierzig sind $3/4$. $\mathcal G$ ist kein Hamiltonoperator; die Differenz $9/4$ ist keine Energielücke. \src{21}

Die Level-1-Vierstromform enthält einen echten Klammerterm zusätzlich zu den Wick-artigen Paarungen. Deshalb können einzelne Stromquadrate verschwinden, während für eine vervollständigte Richtung
\[
 \|J_{-1}(w_{l,s})^2\ket\Omega\|^2=\frac59
\]
gilt. Nach Projektion auf Bell10 beträgt das unnormierte Gewicht $5/18$. Eine explizite positive Frequenzabsorption rekonstruiert die alten Effekte auf einem festgelegten Ein-Strom-Bereich. Eine andere Instrumentierung besitzt jedoch dieselben Effekte und andere Nachzustände. Das physische Mess- und Auffrischgesetz bleibt somit zusätzliche Information. \src{21}

\section{Die tatsächliche Coxeterclock erzwingt mehr als D8}
Im festen CAR-Wörterbuch schickt $C$ die 112 D8-Wurzeln in 56 D8- und 56 Spinorwurzeln. Die kumulative $C$-Bahn erreicht bei Potenz acht alle 240 Wurzelrichtungen. Schon D8 zusammen mit einem geeigneten Spinorstrom und seinem Adjungierten schließt durch nichtverschwindende Klammern in den Schritten $114\to170\to240$. \src{22}

Falls also $C$ auf demselben lokalen Stromraum physisch implementiert ist, kann eine rein D8-beschränkte Theorie nicht unverändert bestehen bleiben. Das ist eine notwendige Schließungsstruktur. Es ist noch kein Bau der 128 fehlenden Spinorstromoperatoren auf dem Rohquellenraum. Auch die gewählte Spinorchiralität wird nicht unabhängig selektiert, wenn $C$ schon aus diesem E8-System stammt.

\section{Der feste komplexe Rahmen ist zu klein}
Die 60 alten Strahlen sind viergliedrige $J$-Orbits vollständiger Wurzeln. Die tatsächliche Coxeterwirkung respektiert diese Quotienten bei festem $J$ nicht: Ein Orbit kann auf zwei Zielorbits aufgeteilt werden. Der korrekte kovariante Ausdruck ist
\[
 ([\alpha]_J,J)\longmapsto
 ([C\alpha]_{CJC^{-1}},CJC^{-1}).
\]
Der letzte Vertrag macht diesen Befund quantitativ. Die Abbildung $A=I+J$ erfüllt $A^TA=2I$ und bildet Standard-E8-Wurzeln der Normquadrate zwei auf die tatsächlichen Gaußschen Quellenwurzeln der Normquadrate vier ab. Sie ist eine skalierte Koordinatenbrücke, kein normgleicher Wurzelautomorphismus. \src{29}

Die spezifizierte Koordinatenbrücke transportiert $C$ zu $C_G=ACA^{-1}$ in den Gaußkoordinaten. Sie wird nicht als einzig mögliche Brückenwahl behauptet. Die Rahmenbahn in diesen Koordinaten
\[
 J_r=C_G^rJC_G^{-r}
\]
hat exakt 15 verschiedene Glieder und entsprechende Ereignisalphabete. Die Projektoren transportieren sich kovariant,
\[
 C_G\Pi_\alpha^{J_r}C_G^{-1}=\Pi_{C_G\alpha}^{J_{r+1}}.
\]
Unter $1\le k\le30$ kommutiert $C_G^k$ nur für $k=15,30$ mit $J$; keine Potenz antikommutiert. Eine zusätzliche Aussage im integralen Standard-Gitterrahmen schließt jede gittererhaltende komplexe Struktur $J_*$ mit $J_*^2=-I$ aus, die mit $C$ kommutiert. Im benutzten Ganzzahlrahmen hat $C$ modulo 31 acht verschiedene Eigenwerte; ein kommutierender Operator müsste auf jeder Eigenlinie skalar wirken, obwohl $-1$ in $\F_{31}$ kein Quadrat ist. Antikommutation scheitert bereits an der ungeraden Halbperiode $C^{15}=-I$. \src{29}

Dieser Satz betrifft den integralen Geltungsbereich, nicht beliebige reelle komplexe Strukturen. Zentralisatoren von Coxeterelementen sind zudem bekannte mathematische Objekte; die Sitzung beansprucht nicht die allgemeine Theorie als Neuerfindung. \cite{coxeter}

Positiv ist eine kovariante Familie $H[J_r]$ möglich. Ihre Existenz wählt aber keine Übergänge zwischen den Rahmen, keine Rahmenenergie und keine physische Eichinterpretation. Die Zahl 15 bezeichnet hier die Größe einer Rahmenbahn, keine neue Zahl von Familien oder Raumdimensionen.

\section{Warum neutrale Paketfluktuationen keine acht Cartanströme tragen}
Auf den Ring-Paketzuordnungen wirkt die Quellen-$J$-Transformation als gemeinsamer Skalar. Sie bleibt nach Grammatrix-Orthonormierung skalar. Auf den tatsächlichen acht reellen Cartanrichtungen gilt dagegen $J^2=-I$, also gibt es keinen Fixvektor. Für eine äquivariante lineare Abbildung $T$ müsste gelten
\[
 T(Jv)=\operatorname{Ad}_{U_J}T(v)=T(v),
\]
und damit $T=0$. Eine mit der Symmetrie verschränkte Dressierung ändert diesen Ausschluss innerhalb des neutralen Raums nicht. \src{29}

Die richtige Lehre lautet nicht, neutrale Observablen zu verwerfen. Ihre Produktstruktur kann über geladene Zwischenzustände vermittelt sein:
\[
 P[A,B]P=[PAP,PBP]+PAQBP-PBQAP,\qquad Q=I-P.
\]
Eine bloße Kompression verliert die letzten beiden Terme. Für die gesuchte Stromantwort müssen deshalb die geladenen Komplementwege mit Energie und spektralem Gewicht mitgeführt werden. Eine passende Anzahl neutraler Zustände reicht nicht.

\section{Was die Clock an der Zeitregel bereits erzwingt}
Unter der ausdrücklichen Prämisse einer phasendiagonalen affinen Stromwirkung besitzt ein Generator die Form
\[
 [H,J_n^\alpha]=(\alpha(h)-an)J_n^\alpha.
\]
Die tatsächliche Coxeterwirkung mit charakteristischem Polynom $\Phi_{30}$ erzwingt $h=0$. Auf dem gemeinsamen irreduziblen Stromraum bleibt somit
\[
 H=aL_0+c.
\]
Die Regel ist bedingt stark: Ist der gemeinsame affine Prozess erst vorhanden, bleibt nur eine positive Zeitkalibrierung. Sie beweist jedoch nicht die bislang vorausgesetzte Invarianz des linearen Stromraums. \src{28}

Eine diskrete Clock und ein gemeinsames Vakuum allein genügen nicht. Gegenbeispiele der Form $H_b=L_0+4bL_0(L_0-1)$ können den geprüften Viertelschritt und den ersten Grad bewahren, aber höhere Grade anders bewegen. Dabei erfüllt jeder Kommutator $\operatorname{ad}_{H_b}$ weiterhin die Leibnizregel. Was scheitert, ist die state-unabhängige skalare Frequenz auf dem linearen Stromraum, nicht die Leibnizregel selbst. \src{28}

\section{Der positive Fermi-Patch: ein kontrollierter Anfang}
Der letzte Herkunftstest arbeitet wieder mit der tatsächlichen CAR-Quelle. Die dortige Einteilchendispersion ist $\varepsilon(k)=-2\cos k$. Auf einem festen, nicht umlaufenden Patch um eine Fermistelle liefert die Skalierung
\[
 e_L(r)=\frac L{2\pi}\sin\frac{2\pi r}L
\]
die Abschätzung
\[
 \bigl|e_L(r)-e_L(r+n)+n\bigr|
 \le\frac{(2\pi)^2}{6L^2}\bigl(|r|^3+|r+n|^3\bigr).
\]
Für passend verschachtelte Träger auf dem festen zusammenhängenden Halbzahl-Patch gilt für die positiven Moden $1,2,3$ zugleich der echte Modenkommutator
\[
 [E^{12}_1,E^{23}_2]=E^{13}_3.
\]
Unter der zusätzlichen gleichen-Spezies-Kinetik $h_8=I_8\otimes h_{\mathrm{v454}}$ gehen die Frequenzen mit kontrolliertem Fehler gegen $-1,-2,-3$; am anderen Fermiast wechseln ihre Vorzeichen. Das ist mehr als eine bloß passende Zustandszählung: Komposition und Zeitfrequenz werden auf demselben ausgeschnittenen CAR-Prozess gemeinsam geprüft. \src{29}

Die Grenze bleibt entscheidend. Acht identische Kinetiken entstehen nur nach der zusätzlichen Wahl $h_8=I_8\otimes h_{\mathrm{v454}}$. Acht Slots erzwingen diese Wahl nicht. Ein fester Patch liefert noch keinen Adjungiertenabschluss, keinen gefüllten Seezustand, keinen Level-1-Zentralterm und keine gemeinsame GNS-Domäne. Die 128 Spinorstromoperatoren fehlen auf demselben Rohraum weiterhin.

Rigorose Gitter-Kontinuum-Konstruktionen für freie Fermionen und WZW-Modelle bieten einen passenden Vergleichsrahmen für die nächsten Konvergenzschritte. Ihre Voraussetzungen müssen aber am TFPT-Quellobjekt nachgewiesen werden. Sie ersetzen weder die Auswahl der Kinetik noch die P1/P2-Identifikation. \cite{osborne}



\chapter{Quellenzeit, lokaler Randlimes und Carry}
% Integration: nach dem Rand-/Gauge-Teil; der QWZ-Befund bleibt quellengebunden.
\section{Einordnung des Strangs}

Die späten Contracts vom 20. September 2026 verfolgen eine einzige präzise
Frage: Lässt sich der algebraisch bereits formulierte E8-/Spinor-Anschluss als
derselbe physische Prozess wiederfinden, der Zustand, geladene Felder, Zeit,
Cocycles und Auslesung der ursprünglichen Quelle gemeinsam trägt? Die
Antwort der geprüften Texte ist partiell, aber deutlich. Mehrere starke
Teilresultate sind inzwischen exakt formuliert: Der unveränderte QWZ-Vertreter
hat in der festbreiten, fest skalierten Quelle nur die bekannten Randzweige;
auf einer gemeinsam wachsenden Kreisgeometrie konvergiert ein lokales geladenes
CAR-Feld samt seiner Zeitentwicklung zu einem einzelnen chiralen freien Feld;
eine ausdrücklich vorgegebene Fluxfahrt transportiert im vollen See eine
regionale Ladung. Auf der Zielseite sind der geladene Strompfad, die vier
Diskriminantsektoren, die native Paarstruktur und ihre zeitabhängigen
Antworten algebraisch bestimmt.

Diese Resultate bleiben an ihre Bedingungen gebunden. Die Querbreite eines
QWZ-Streifens zählt keine unabhängigen Spezies, ein E8-Gitter ist keine aus
P1/P2 abgeleitete Quelle, und ein positiver oder unitärer Hilfsanschluss ist
keine Herleitung der nativen Wechselwirkung. Die Contracts tragen deshalb
durchgehend den Status \texttt{\detokenize{PARTIAL}}. Es wurde kein T1--T8-Gate, keine vollständige
TOE und kein 3+1-dimensionaler Ursprung geschlossen.

Die Vertragskennungen und Beweisdateien stehen im Quellenregister.

\section{Was die unveränderte QWZ-Quelle tatsächlich enthält}

Der Contract \src{31} analysiert die tatsächlichen
Hopping-Blöcke aus \path{verification/v1033_charged_disorder.py}, mit Masse eins,
fester Breite $W=8$, Sektor $r=1$ und dem gefüllten negativen See. Für

\[
 p_j=\frac{2\pi(j-1/4)}{N},\qquad D_N=\frac{N}{2\pi}h_N,\qquad
 H_N=d\Gamma(D_N)-\operatorname{Tr}(1_{D_N<0}D_N),
\]

zerfällt der Einteilchenblock in eine obere und eine untere Kante sowie
massive Querzweige. Mit \(\rho=1-\cos p\) und dem unteren Shift \(S_W\) gilt

\[
 B=\rho I-S_W,
\qquad B^*B=T-|e_\mathrm{last}\rangle\langle e_\mathrm{last}|,
\quad T=(1+\rho^2)I-\rho(S_W+S_W^*).
\]

Rang-eins-Interlacing liefert für die nichtkleinen Singulärwerte eine
uniforme Schranke. Für die $2W-2$ Bulkenergien folgt

\[
 |e_\mathrm{bulk}(p)|\ge
 \sqrt{\sin^2p+(1-\rho)^2}=1.
\]

Die übrigen beiden Zweige können nur für \(|p|<\pi/2\) unterhalb dieser
Schranke liegen. Die Fockraum-Ausschöpfung ist dabei nicht nur eine
Einteilchenbehauptung: Für jedes feste \(E<\infty\) und hinreichend großes \(N\) liegen alle Zustände in \(1_{[0,E]}(H_N)\) im endlichen Impulsfenster der beiden Kantenzweige. Teilchen- und Lochanregungen werden gemeinsam
erfasst; Bulkzustände können die Energie nicht durch eine unzulässige
Teilchen-Loch-Kompensation absenken.

Das ist ein positives Resultat über diese ausgewählte Quelle und ihre
Zeitskalierung. Es schließt in dieser Darstellung zusätzliche verborgene
gebundene Bulkkanäle aus. Es beweist weder, dass der rohe P1/P2-Kern genau
diese QWZ-Quelle ist, noch eine lokale Scaling-Algebra-Ausschöpfung für alle
Operatorprodukte oder wechselnde Breiten, Fluxsektoren und Wechselwirkungen.
Insbesondere hat ein unveränderter Rand mit einem komplexen Zweig pro
Chiralität nicht automatisch acht E8-Kanäle. Die Aussage „acht Querreihen"
wäre eine Verwechslung von Ortsreihen mit Spezies.

Der gleiche Contract zeigt eine zweite Grenze. Die abstrakte E8-Theorie hat
bei Grad eins $1+248=249$ Zustände, während der QWZ-Grenzraum unterhalb der
festgelegten Quelleinheiten bei \(E<1{,}1\) nur Rang zehn besitzt. Dieser
Vergleich setzt die dortige Energieeichung voraus und ist keine universelle
Unmöglichkeitsaussage gegen P1/P2. Er markiert aber, dass E8-Grade und rohe
Quellenergie nicht ohne gemeinsame Kalibrierung identifiziert werden dürfen.

Die Fixpunkt-Abkürzung löst diese Frage ebenfalls nicht. Für einen
Punktdefekt ist die Kopplungsdimension $1-h$; für eine ausgedehnte
1+1-dimensionale Kante ist sie $2-(h+\bar h)$. Ein chirales Objekt mit
$(h,\bar h)=(2,0)$ ist dort nach Potenzzählung marginal. Die Contracts
verwenden dies nicht als Beta-Funktionsbeweis, sondern als Ausschluss der
Behauptung, Potenzzählung allein wähle eine eindeutige Kantendynamik.

\section{Der lokale wachsende-Kreis-Limes}

\src{33} prüft eine andere Geometrie als die früheren
Festkreisvergleiche: \(R\to\infty\) und \(a=2\pi R/N\to0\), bei unveränderter
Masse, Breite und Hoppingstruktur. Für den tatsächlichen Streifenoperator und
den oberen Randvektor \(r_\mathrm{top}\) gilt exakt

\[
 h(p)r_\mathrm{top}=-\sin(p)r_\mathrm{top}+(1-\cos p)\chi_\mathrm{top}.
\]

Daraus ergibt sich für \(D_a(k)=h(ak)/a\)

\[
 \|(D_a(k)+k)r_\mathrm{top}\|
 \le \varepsilon_a(k)
 =\frac{a k^2}{2}+\frac{a^2|k|^3}{6}.
\]

Die Spektralprojektion auf den gefüllten negativen See und die echte
Zeitentwicklung werden für $k\ne0$ gleichzeitig kontrolliert:

\[
 \left|\langle r_\mathrm{top},P_a e^{-itD_a}r_\mathrm{top}\rangle
 -1_{k>0}e^{itk}\right|
 \le |t|\varepsilon_a(k)+\frac{\varepsilon_a(k)^2}{k^2}.
\]

Auf jedem festen \(k\)-Band und beschränkten Zeitintervall verschwindet die
rechte Seite. Nach Fourier-Sampling mit

\[
 (I_{R,a}f)_j=R^{-1/2}\widehat f(k_j)r_\mathrm{top}
\]

konvergieren die CAR-Kerne, die gefüllte Kovarianz und alle festen endlichen
Wörter aus Feldern, Adjungierten und Zeitübersetzungen. Der Grenzgenerator ist
\(i\partial_x\); es entsteht ein einzelnes chirales freies Fermion. Die
Viertelholonomie wird im festgehaltenen lokalen Beobachtungsbereich unsichtbar, während der physische Kreisumfang wächst. Das ist ein echter positiver Quellen-/Zeitbefund,
kein Fit an eine isolierte Dispersion.

Der Contract misst zugleich die Kanalzahl am selben Quelloperator: Bei
$p=0$ hat der Block genau zwei Nullmoden, eine pro Rand, und vierzehn
massive Zweige. Die lokale Niedrigenergieskalierung besitzt also einen
chiralen Kanal je Kante mit entgegengesetzter Geschwindigkeit. Der positive
Limes beseitigt die frühere Festkreisabweichung nicht durch bloße
Gitterverfeinerung: Bei festem \(R\) bleibt die Holonomie-Differenz bestehen.
Die entscheidende Zusatzannahme ist der wachsende physische Umfang.

Damit sind drei Ebenen sauber getrennt: ein lokaler Ein-Kanal-Limes der
ausgewählten QWZ-Quelle, ein bedingtes achtkanaliges E8-Ziel und die noch
fehlende Abbildung des einen in das andere. Der Limes liefert weder acht
Spezies noch den Halbspinor, die E8-Cocycle oder eine P1/P2-Herkunft.

\section{Zeit, geladene Ströme und Viertelholonomie}

\src{32} verbindet die vorhandene Gittervertexalgebra mit
der Frage nach gemeinsamer Quellzeit. Der bekannte E8-Pfad ist nicht null:
mit \(\alpha=(-1,-1,0,\ldots)\), \(\beta=s=(1/2)^8\) und
\(\gamma=\alpha+s\) gilt \(\alpha\cdot s=-1\), \(\lvert\gamma\rvert^2=2\), und der
einfache OPE-Pol führt von \(V_\alpha\) und \(V_s\) nach \(V_\gamma\).
Die einzelne Phase hängt von der Cocycle-Konvention ab; der nichtverschwindende
Pfad und die relativen Kommutatorzeichen sind invariant. Für das elementare
Vektorfeld \(e_1\) ist \(e_1\cdot s=1/2\), also die Monodromie (-1). Das
zeigt, dass nacktes Vektorfeld und Spinorstrom in der bosonischen E8-
Vakuumalgebra unterschiedliche Lokalitätsrollen haben.

Die CAR-Stromgrenze ist im endlichen Fenster exakt, sofern die Spezies
verschieden sind. Für

\[
 E_n^{ab}=\sum_r b^*_{a,r}b_{b,r-n}
\]

gilt auf dem stabilisierten Kern

\[
 [E_m^{ab},E_n^{cd}]
 =\delta_{bc}E_{m+n}^{ad}-\delta_{ad}E_{m+n}^{cb}
 +m\delta_{m+n,0}\delta_{ad}\delta_{bc}I.
\]

Für gleiche Spezies wäre die analoge Paarformel falsch, weil
\(A_n^{aa}=0\) durch Antisymmetrie, während die behauptete rechte Seite auf
dem Vakuum nicht verschwindet. Dieser Speziesvorbehalt ist eine konkrete
Lücke in der Paarübertragung.

Die rohe Viertelholonomie wird durch die zeitabhängige Ladung sichtbar. Die
Quelle trägt

\[
 E_\mathrm{source}(q)=q^2/2-q/4,
 \qquad H_\mathrm{source}=L_0-Q/4,
\]

während das bedingte achtkanalige Ziel \(H_\delta=L_0-\delta\cdot P\) mit
\(\delta=$1/4$^8=s/2\) verwendet. Für einen geladenen Vertex gilt

\[
 [H_\delta,V_{p,n}]=(-n-\delta\cdot p)V_{p,n}.
\]

Zahlenerhaltende Differenzen \(p=e_a-e_b\) sehen den Zusatzterm nicht. Ein
Paarstrom erfährt dagegen (-1/2), der Halbspinor (p=s) (-1), und der
entgegengesetzte Spinor (+1). Deshalb kann ein neutraler Stromtest bestehen,
obwohl die gemeinsame geladene Zeit falsch identifiziert ist. Dasselbe gilt
für ein formal normiertes Feld \(V_{s,-1}\Omega\): sein konformes Gewicht ist
eins, seine \(H_\delta\)-Energie bei Viertelholonomie aber null.

Die rationalen Clock-Lifts (C,J) können auf dem E8-Gitter mit Cocycle
ausgeführt werden; ihre Ordnungen bleiben erhalten. Unter der festen Zeit

\[
 W_g(t)=e^{itH_\delta}U_g e^{-itH_\delta}
 =e^{it(g\delta-\delta)\cdot P}U_g.
\]

Das ist ein exakter Heisenbergtransport und keine erhaltene Symmetrie. Die
Ladungsphasen müssen in einem gemeinsamen Quellenvergleich auftauchen; sie
werden nicht durch Umbenennung entfernt. Innerhalb der vorgegebenen Klasse
homogener positiver quadratischer Stromhamiltonoperatoren mit linearer
Ladungskorrektur erzwingt gleichzeitige \(C\)- und \(J\)-Invarianz
(V=vI) und (b=0). Diese Auswahl ist jedoch nur dann physisch relevant,
wenn die ursprüngliche Quelle beide Clocks als Symmetrien derselben Zeit
realisiert.

\section{Carry: positive Quelle, aber externe Fahrregel}

\src{36} bestätigt die Quellformeln, markiert aber die
Fensterbedingung: Die Stromgewichte gelten für die deklarierte
Niedrigimpulsmenge, nicht als ungefilterter Vollquellensatz. Die E8-Cocycle
trägt den Carry \(T_s^2=\epsilon(s,s)T_{2s}\) mit
\(\epsilon(s,s)=-1\) in der übernommenen Konvention. Ein phasenfreier Lift
und der echte Cocycle unterscheiden sich nur in der Rephasierung, nicht in
der Tatsache des Carry.

Der stärkere Contract \src{35} setzt dann ausdrücklich
eine externe Fluxfahrt voraus. Für $W=8$, Masse eins und \(N\ge32\) wird

\[
 f(s)=1/4+s,\qquad i\partial_sU(s)=T h_N(f(s))U(s),\qquad U(0)=I
\]

vorgegeben. Der volle negative See wird ohne Reset, Bad oder Nachbesetzung
entwickelt; die Teilchenzahl bleibt \(8N\). Genau am Nullübergang werden die
beiden Kantenlinien gemeinsam verfolgt. Für einen geeigneten Kato-
Intertwiner liefert der Contract

\[
 \|\Gamma-P_\mathrm{tr}\|_1=O(T^{-1}),\qquad
 \|\Gamma-P_\mathrm{tr}\|_2=O((T\sqrt N)^{-1}).
\]

Mit einem gaußgefilterten regionalen Ladungsoperator erhält man

\[
 \langle C_\mathrm{top}(3/4)\rangle
 =1/2+O(N^{-1})+O(T^{-1}),
\]

und am vollen Fluxzyklus

\[
 \langle C_\mathrm{top}(5/4)\rangle
 =1+O(N^{-16})+O(T^{-1}).
\]

Für \(T=N^{3/4}\) werden zugleich Ladungs-, Zustands- und Energiegraph-
Fehler kontrolliert; der feste Randphasenfehler ist $O(N^{-1/4})$. Das ist
ein positiver, quellennaher Carry im vollen Fockraum. Die Operation wirkt
jedoch auf allen Umfangslinks, ihre Dauer divergiert in Hoppingeinheiten,
und bei zeitabhängigen Eichwechseln ist der Term \(i\dot G G^*\) nötig. Sie
ist daher kein lokaler Seam-Vertex und keine automatisch lokale Spinorfeld-
Einfügung. Der Contract sagt ausdrücklich: Ein globaler Spectral-flow-
Operator liefert weder lokale Produkte noch die achtkanalige Cocycle.


